\begin{example}
    La sucesión $\partfrac{n^2\alpha}$ con $\alpha \notin \mathbb{Q}$ es equidistribuida en $[0,1)$.
\end{example}

\begin{resolve}
    Sea $f(n)=n^2\alpha$ con $\alpha \notin \mathbb{Q}$, entonces
    \[
        f(n+h)-f(n)=n^2\alpha + 2nh\alpha + h^2\alpha - n^2\alpha = 2nh\alpha+h^2\alpha.
    \]

    Usando la estimación anterior, como para $h=0$ se tiene $e^{2\pi i 0}=1$, entonces
    \[
        \left|\sum_{n=1}^N e^{2\pi i  k n^2 \alpha}\right|^2 \leq 
        4\,\frac{N^2}{H} + 4\,\frac{N}{H} \sum_{h=1}^{H-1} \left|\sum_{n=1}^{N-h} e^{2\pi i k (2nh\alpha+h^2\alpha)}\right|.
    \]

    Para cada $h \neq 0$, por el ejemplo \ref{ex:ax-b}, se tiene que $\partfrac{2nh\alpha+h^2\alpha}$ es equidistribuida en $[0,1)$.

    Al dividir por $N^2$,
    \[
        \left|\frac{1}{N}\sum_{n=1}^N e^{2\pi i n^2 \alpha}\right|^2 \leq 
        \frac{4}{H} + \frac{4}{N} \sum_{h=1}^{H-1} \frac{N-h}{N} \left|\frac{1}{N-h}\sum_{n=1}^{N-h} e^{2\pi i k(2nh\alpha+h^2\alpha)}\right|.
    \]

    Fijado $H$, por el criterio de Weyl, para cada $h \neq 0$ se tiene que
    \[
        \lim_{N\to\infty} \frac{1}{N-h}\sum_{n=1}^{N-h} e^{2\pi i k(2nh\alpha+h^2\alpha)} = 0.
    \] 

    Luego,
    \[
        \limsup_{N\to\infty} \left| \frac{1}{N} \sum_{n=1}^N e^{2\pi i k n^2 \alpha}\right|^2 \leq \frac{4}{H}.
    \]  

    Dado que es cierto para cada $H\in \mathbb{N}$, haciendo $H\to\infty$ se concluye que
    \[
        \lim_{N\to\infty} \left| \frac{1}{N} \sum_{n=1}^N e^{2\pi i k n^2 \alpha}\right|^2 = 0.
    \]
    
    Por lo tanto, por el criterio de Weyl, la sucesión $\partfrac{n^2\alpha}$ es equidistribuida en $[0,1)$.
\end{resolve}